\section{Proof of the main theorem}\label{sec:assembly}

Fix $f$ as in \cref{thm:main}. The finiteness and positivity of $C_f$ are \cref{prop:poly}(0). Let $x\ge2$ and $y\ge1$ be real with $Y=f(y)\le x$, and fix $\eps>0$.

\emph{Range of the moduli.} The hypothesis $f'(t)\ge1$ for all \emph{real} $t\ge1$ makes $f$ increasing on $[1,\infty)$. Hence every integer $1\le k\le y$ satisfies
\begin{equation}\label{eq:moduli-range}
2\le f(k)\le f(y)=Y\le x .
\end{equation}
So \cref{prop:single} applies to $q=f(k)$ with the same $x$, for every $k\le y$, with a constant independent of $k$. By~\eqref{eq:Vf}, $V_f(x,y)=\sum_{1\le k\le y}f'(k)V(x;f(k))$, and $f'(k)>0$.

\emph{Elementary inequalities.} Write $f(t)=\sum_{i\le d}\alpha_it^i$ and $\gamma_0=\sum_i|\alpha_i|$. For $y\ge1$, $Y\le\gamma_0y^d$, hence
\begin{equation}\label{eq:y-Y}
y^{-1/2}\le\gamma_0^{1/(2d)}\,Y^{-1/(2d)}.
\end{equation}
Since $Y\le x$, we have $Y^\eps\le x^\eps$ and $1+\log Y\le1+\log x\le C'(\eps)x^\eps$. With \cref{prop:poly}(b),(c) this gives
\begin{equation}\label{eq:sums}
\sum_{k\le y}\frac{f'(k)}{f(k)}\le c_8(f)C'(\eps)\,x^{\eps},\qquad\sum_{k\le y}f'(k)\le c_8(f)\,Y .
\end{equation}

\subsection{Part (ii)}

Apply \cref{prop:single}(i) with $\eps/2$ to each $q=f(k)$, multiply by $f'(k)$ and sum over $1\le k\le y$:
\[
\Bigl|V_f(x,y)-\sum_{k\le y}f'(k)\MT(x;f(k))\Bigr|\le c_7(\tfrac\eps2)\Bigl(x^{5/3+\eps/2}\sum_{k\le y}\frac{f'(k)}{f(k)}+\sqrt x\sum_{k\le y}f'(k)\Bigr).
\]
By~\eqref{eq:sums} with $\eps/2$, the right-hand side is at most
\[
c_7(\tfrac\eps2)\,c_8(f)\bigl(C'(\tfrac\eps2)\,x^{5/3+\eps}+x^{1/2}Y\bigr),
\]
which is part~(ii).

\subsection{Part (i)}

Apply \cref{prop:single}(ii) with $\eps/3$ to each $q=f(k)$ and use $f'(k)\sqrt{f(k)}=F(k)$:
\[
\Bigl|V_f(x,y)-\sqrt x\sum_{k\le y}F(k)h(f(k))\Bigr|\le c_7(\tfrac\eps3)\,(T_1+T_2+T_3),
\]
where the three sums are bounded as follows.
\begin{itemize}
\item $T_1=x^{1/3}\sum_{k\le y}f'(k)f(k)^{2/3+\eps/3}\le x^{1/3}Y^{2/3+\eps/3}\sum_{k\le y}f'(k)\le c_8(f)\,x^{1/3+\eps/3}Y^{5/3}$, by $f(k)\le Y\le x$ and~\eqref{eq:sums}.
\item $T_2=x^{5/3+\eps/3}\sum_{k\le y}f'(k)/f(k)\le c_8(f)C'(\tfrac\eps3)\,x^{5/3+2\eps/3}$, by~\eqref{eq:sums} with $\eps/3$.
\item $T_3=\sqrt x\sum_{k\le y}f'(k)\le c_8(f)\,x^{1/2}Y\le c_8(f)\,x^{1/2}Y^{3/2-1/(2d)}$, because $Y\ge2$ and $\frac12-\frac1{2d}\ge0$.
\end{itemize}
For the main term, \cref{prop:poly}(a) with $\eps/3$ and~\eqref{eq:y-Y} give
\begin{align*}
\Bigl|\sqrt x\sum_{k\le y}F(k)h(f(k))-C_fx^{1/2}Y^{3/2}\Bigr|&\le c_8(f,\tfrac\eps3)\gamma_0^{1/(2d)}\,x^{1/2}Y^{3/2+\eps/3-1/(2d)}\\
&\le c_8(f,\tfrac\eps3)\gamma_0^{1/(2d)}\,x^{1/2+\eps/3}Y^{3/2-1/(2d)}.
\end{align*}
By the triangle inequality,
\[
\bigl|V_f(x,y)-C_fx^{1/2}Y^{3/2}\bigr|\le c(f,\eps)\Bigl(x^{1/3+\eps}Y^{5/3}+x^{5/3+\eps}+x^{1/2+\eps}Y^{3/2-1/(2d)}\Bigr),
\]
with $c(f,\eps)=c_7(\tfrac\eps3)c_8(f)\bigl(2+C'(\tfrac\eps3)\bigr)+c_8(f,\tfrac\eps3)\gamma_0^{1/(2d)}$. This is part~(i).

\subsection{Part (iii)}

Fix $\delta\in(0,1/9)$. The interval $[x^{7/9+\delta},x^{1-\delta}]$ is nonempty for $x>1$ exactly because $7/9+\delta<1-\delta$. Assume $x^{7/9+\delta}\le Y\le x^{1-\delta}$, put $\eta=\min(\delta/12,1/(6d))$, and apply~(i) with $\eps=\min(\delta/12,1/(9d))$. We divide each error term by $x^{1/2}Y^{3/2}$:
\begin{enumerate}[label=\arabic*.]
\item $x^{-1/6+\eps}Y^{1/6}\le x^{-1/6+\eps+(1-\delta)/6}=x^{\eps-\delta/6}\le x^{-\delta/12}\le x^{-\eta}$, using $Y\le x^{1-\delta}$ and $\eps\le\delta/12$;
\item $x^{7/6+\eps}Y^{-3/2}\le x^{7/6+\eps-\frac32(7/9+\delta)}=x^{\eps-3\delta/2}\le x^{-\delta/12}\le x^{-\eta}$, using $Y\ge x^{7/9+\delta}$ and $\frac32\cdot\frac79=\frac76$;
\item $x^{\eps}Y^{-1/(2d)}\le x^{\eps-7/(18d)}\le x^{1/(9d)-7/(18d)}=x^{-5/(18d)}\le x^{-1/(6d)}\le x^{-\eta}$, using $Y\ge x^{7/9}$ and $\eps\le1/(9d)$.
\end{enumerate}
Hence $|V_f(x,y)-C_fx^{1/2}Y^{3/2}|\le3c(f,\eps)\,x^{-\eta}x^{1/2}Y^{3/2}$, where $\eps$ depends only on $\delta$ and $d$. As $0<C_f<\infty$, this is $V_f(x,y)=C_fx^{1/2}Y^{3/2}\bigl(1+O_{f,\delta}(x^{-\eta})\bigr)$. \qed

\begin{remark}
The exponent $\eta$ is conservative. The third ratio is in fact at most $x^{-5/(18d)}$, and the second at most $x^{-17\delta/12}$. We keep the exponent as it was stated in the archived proof.
\end{remark}
